\scanpage{450}
\mspage{362}
\inserted{on peut $=$ \unsure{prendre|préciser}} $K'$ telle que $\mathrm{Ob}\,K'=\mathrm{Ob}\,K\unsure{\amalg|\sqcup|\cup}\{e\}$, où $e$ est un \uline{objet final} de $K'$, et $K$ \deleted{\uline{Dém}} \deleted{\unsure{étant}} \inserted{est} une sous-catégorie pleine de $K'$. (Donc \inserted{si $K'\ne K$,} $K'$ est isom. à un cône à droite construit sur $K$, ``déduit de $K$ en lui \unsure{ajoutant} un objet final''.)

\inserted{\uline{Dém} :} Si $K$ admet un objet final, on prend $K'=K$. Sinon, Comme $J$ est grde devant $\mathrm{Card}\,\mathrm{Ob}\,K$, il existe un $j\in\mathrm{Ob}\,J$ qui majore tous les objets de $K$. \deleted{\unsure{Choisissons pour}} \inserted{Choisissons}, \deleted{Pour tout $k\in\mathrm{Ob}\,K$, considérons alors flèches \unsure{toutes les}} pour tt $k\in\mathrm{Ob}\,K$, une flèche
\[
\alpha_k:k\longrightarrow j.
\]
Pour tte flèche
\[
u:k'\unsure{\text{---}|\longrightarrow}k
\]
de $K$, considérons la double-flèche
\[
(\alpha_k\circ u,\alpha_{\unsure{k'|m}}):k'\rightrightarrows j.
\]
L'ens de ces doubles flèches, de \deleted{\unsure{but}} $j$, est de \unsure{cardinal|cardinalité} $\leq\mathrm{Card}\,\mathrm{Fl}\,K\leq\pi$, donc $J$ étant grde devant $\pi$ il existe une flèche $\beta:j\unsure{\text{---}|\longrightarrow}e$ de $J$ qui égalise ttes ces doubles-flèches. Soit $K'\subset J$ ayant comme objets $\mathrm{Ob}\,K\cup\{e\}$, et comme \deleted{\unsure{couples}} flèches, \deleted{en plus} \unsure{En} plus de
% archive sheet number: 449
